% !TEX root = ../thesis.tex
\chapter{Cơ sở lý thuyết}

Chương này xây dựng nền tảng lý thuyết cho toàn bộ các mô hình và thuật toán được sử dụng trong HLK-ViT5.
Mục~2.1 trình bày kiến trúc Transformer Encoder-Decoder và mô hình ViT5-base;
Mục~2.2 trình bày các kỹ thuật nén ngữ cảnh và trích chọn thông tin như MMR và Keyword Conditioning;
Mục~2.3 giới thiệu các mô hình baseline đối chứng;
Mục~2.4 định nghĩa chi tiết 4 nhóm độ đo đánh giá chất lượng tóm tắt; và Mục~2.5 tổng kết chương.

\section{Mô hình ViT5 và Kiến trúc Transformer}

\subsection{Kiến trúc Transformer Encoder-Decoder}

Kiến trúc Transformer \cite{Vaswani2017} dựa trên cơ chế chú ý (Attention Mechanism), loại bỏ hoàn toàn các vòng lặp hồi quy (RNN) để cho phép tính toán song song hóa trên quy mô lớn. Đối với bài toán sinh văn bản dạng Sequence-to-Sequence, kiến trúc Encoder-Decoder bao gồm hai khối chính:

\begin{itemize}
  \item \textbf{Bộ mã hóa (Encoder):} Tiếp nhận chuỗi token đầu vào $X = (x_1, x_2, \ldots, x_n)$, sử dụng các lớp chú ý bản thân đa đầu (Multi-Head Self-Attention) và mạng truyền thẳng (Feed-Forward Neural Network) để tạo ra tập hợp các biểu diễn ngữ cảnh ẩn $H = (h_1, h_2, \ldots, h_n)$.
  \item \textbf{Bộ giải mã (Decoder):} Tự điều khiển sinh chuỗi đầu ra $Y = (y_1, y_2, \ldots, y_m)$ theo từng bước (autoregressive). Tại mỗi bước $t$, bộ giải mã kết hợp chú ý bản thân có mặt nạ (Masked Self-Attention) trên các token đã sinh $(y_1, \ldots, y_{t-1})$ và chú ý chéo (Cross-Attention) trên chuỗi ẩn $H$ của Encoder.
\end{itemize}

Công thức chú ý tích vô hướng có tỷ lệ (Scaled Dot-Product Attention) với truy vấn $Q$, khóa $K$ và giá trị $V$ được định nghĩa:

\[
\text{Attention}(Q,K,V) \;=\; \text{softmax}\!\left(\frac{QK^{\top}}{\sqrt{d_k}}\right) V
\]

\noindent trong đó $d_k$ là chiều của khóa.

\subsection{Mô hình tiền huấn luyện ViT5}

ViT5 \cite{Nguyen2022} là mô hình ngôn ngữ tiếng Việt được xây dựng dựa trên kiến trúc T5 (Text-to-Text Transfer Transformer). ViT5 chuyển đổi tất cả các tác vụ NLP (phân loại, dịch máy, tóm tắt) về dạng bài toán chuỗi-sang-chuỗi: nhận đầu vào là một chuỗi văn bản và sinh đầu ra là một chuỗi văn bản.

Phiên bản \texttt{VietAI/ViT5-base} có 12 lớp Encoder, 12 lớp Decoder, chiều ẩn $d_{\text{model}} = 768$, 12 đầu chú ý và khoảng 220 triệu tham số. Mô hình được tiền huấn luyện trên tập dữ liệu tiếng Việt quy mô lớn với mục tiêu khôi phục token bị che (Span Corruption). ViT5 là mô hình nền tảng được chọn để xây dựng và tinh chỉnh trong đề tài HLK-ViT5.

\section{Các kỹ thuật xử lý ngữ cảnh và trích chọn thông tin}

\subsection{Sentence-aware Chunking và Overlapping}

Khi độ dài bài báo vượt quá cửa sổ xử lý của Encoder, văn bản cần được phân rã thành các đoạn nhỏ (chunks). Kỹ thuật phân đoạn theo ranh giới câu (\textit{Sentence-aware Chunking}) đảm bảo không cắt đôi câu giữa chừng, giữ nguyên tính hoàn chỉnh cú pháp. Thêm vào đó, việc duy trì khoảng đè 1 câu (\textit{1-sentence Overlapping}) giữa hai chunk kề nhau giúp bảo toàn mối quan hệ tham chiếu ngữ nghĩa giữa các đoạn.

\subsection{Maximal Marginal Relevance (MMR)}

Thuật toán Maximal Marginal Relevance (MMR) được thiết kế nhằm chọn ra tập hợp các câu/đoạn văn có độ tương quan cao nhất với chủ đề nhưng đồng thời tối thiểu hóa sự trùng lặp thông tin nội bộ. Với tập hợp ứng viên $C$, tập đã chọn $S$, truy vấn/từ khóa $Q$, chỉ số MMR của một ứng viên $c_i \in C \setminus S$ được tính:

\[
\text{MMR}(c_i) \;=\; \arg\max_{c_i \in C \setminus S} \left[ \lambda \cdot \text{Sim}_1(c_i, Q) - (1 - \lambda) \cdot \max_{c_j \in S} \text{Sim}_2(c_i, c_j) \right]
\]

\noindent trong đó $\text{Sim}_1$ đo độ nổi bật của chunk so với bài báo/từ khóa, $\text{Sim}_2$ đo độ tương đồng giữa hai chunk, và $\lambda \in [0,1]$ là hệ số cân bằng giữa độ nổi bật và tính đa dạng.

\subsection{Keyword-conditioned Generation}

Keyword-conditioned Generation là kỹ thuật bổ sung danh sách từ khóa đại diện vào trước ngữ cảnh đầu vào của bộ mã hóa. Bằng cách định dạng chuỗi đầu vào dưới dạng:

\begin{center}
\texttt{Keywords: [từ\_khóa\_1, từ\_khóa\_2, ...] | Context: [nội\_dung\_đã\_lọc]}
\end{center}

\noindent mô hình ViT5 được gợi ý tập trung chú ý vào các thực thể và khái niệm trung tâm khi sinh bản tóm tắt.

\section{Các mô hình đối chứng (Baselines)}

Đề tài đưa hai mô hình mã nguồn mở hiện có vào so sánh đối chứng:

\begin{itemize}
  \item \textbf{\texttt{Nishikyen/vit5-vietnamese-news}:} Mô hình ViT5-base được fine-tune trên tập dữ liệu tin tức tiếng Việt theo phương pháp chuẩn hóa truyền thống với chiều dài input giới hạn 256/512 tokens.
  \item \textbf{\texttt{thnhan/sft\_model}:} Mô hình ViT5-base được tinh chỉnh theo phương pháp SFT (Supervised Fine-Tuning) tập trung vào khả năng sinh câu trôi chảy.
\end{itemize}

\section{Hệ đo lường đánh giá chất lượng tóm tắt}

Để đánh giá toàn diện HLK-ViT5, nghiên cứu xây dựng bộ đo lường 4 nhóm:

\subsection{Nhóm 1: Lexical Overlap (Độ trùng lặp từ vựng)}

\begin{itemize}
  \item \textbf{ROUGE-1 / ROUGE-2:} Đo độ trùng lặp $n$-gram ($n=1, 2$) giữa bản tóm tắt mô hình $C$ và tham chiếu $R$.
  \item \textbf{ROUGE-L:} Đo độ dài chuỗi con chung dài nhất (Longest Common Subsequence).
\end{itemize}

\subsection{Nhóm 2: Semantic Quality (Chất lượng ngữ nghĩa)}

\begin{itemize}
  \item \textbf{BERTScore (Precision, Recall, F1):} Sử dụng vector nhúng ngữ cảnh từ mô hình PhoBERT \cite{Nguyen2020} để tính độ tương đồng cosine tốt nhất giữa các token, phản ánh sự tương đồng ngữ nghĩa ngay cả khi từ ngữ được diễn đạt lại.
\end{itemize}

\subsection{Nhóm 3: Factual Consistency (Tính nhất quán sự thật)}

\begin{itemize}
  \item \textbf{Entity Precision:} Tỷ lệ phần trăm các thực thể (tên riêng, tổ chức, địa danh) xuất hiện trong bản tóm tắt được xác nhận có có mặt chính xác trong bài gốc.
  \item \textbf{Number Consistency:} Tỷ lệ bảo toàn chính xác các con số và dữ liệu thống kê.
  \item \textbf{Contradiction Rate / Hallucination Rate:} Tỷ lệ phần trăm các mệnh đề trong bản tóm tắt mâu thuẫn hoặc "bịa" thêm so với nguồn gốc.
\end{itemize}

\subsection{Nhóm 4: Performance (Hiệu năng tính toán)}

\begin{itemize}
  \item \textbf{Inference Time:} Thời gian xử lý trung bình (tính bằng giây) để tạo ra bản tóm tắt hoàn chỉnh cho một văn bản đầu vào.
\end{itemize}

\section{Tổng kết chương}

Chương này đã cung cấp nền tảng lý thuyết về mô hình ViT5, kiến trúc Transformer Encoder-Decoder, thuật toán MMR, kỹ thuật Keyword Conditioning và hệ đo lường 4 nhóm. Chương~3 sẽ trình bày cách các lý thuyết này được tích hợp vào kiến trúc pipeline HLK-ViT5.
